\section*{Capítulo \ref{ch:b1:radiation-heat-transport} --- Radiación y transporte de calor}
% values: calc/ch27.py

\begin{solution}{exo:b1:radiation-heat-transport:1}
(a)~$S_0/\pi = 1361/\pi = \qty{433.2}{W.m^{-2}}$. (b)~En el ecuador, en un equinoccio,
$\bar Q = (S_0/\pi)(a/r)^2$, así que $(a/r)^2 = 436.7/433.2 = 1.0080$ y $r/a =
1.0080^{-1/2} = 0.9960$: la Tierra está un \qty{0.4}{\percent} más cerca del Sol que su distancia media. (c)~$\bar Q = 433.2 \times 0.9683 \times \cos\ang{23.44} =
\qty{384.9}{W.m^{-2}}$.
\end{solution}

\begin{solution}{exo:b1:radiation-heat-transport:2}
(a)~$415.6/172.4 = 2.41$. (b)~La zona entre los paralelos $\pm\varphi$ tiene
el área $2\pi a_\oplus^2 \times 2\sin\varphi$, una fracción $\sin\varphi$ de la esfera:
$\sin\ang{37.0} = 0.602$, \qty{60}{\percent}. (c)~En un eje de \omterm{def:b1:spherical-astronomy:geographic}{latitud}, las bandas cerca de los polos parecen tan anchas como las cercanas al ecuador, pero cubren mucha menos área: la zona hacia el polo desde \ang{60} es solo $1 - \sin\ang{60} =
\qty{13.4}{\percent}$ de la superficie. Frente a $\sin\varphi$, los trópicos llenan la mayor parte del eje.
\end{solution}

\begin{solution}{exo:b1:radiation-heat-transport:3}
(a)~$T_e = (\text{OLR}/\sigma)^{1/4}$: \qty{256.4}{K}, \qty{241.3}{K} y
\qty{254.5}{K}. (b)~Luz solar absorbida: $287.7/75.8 = 3.80$; OLR: $244.9/192.3 =
1.27$. (c)~La \omterm{def:b1:radiation-heat-transport:net-radiation}{radiación neta} sigue a la luz solar absorbida: positiva y grande en los trópicos, fuertemente negativa en las altas \omterm{def:b1:spherical-astronomy:geographic}{latitudes}. La OLR, que varía poco, no compensa la caída de la luz solar absorbida.
\end{solution}

\begin{solution}{exo:b1:radiation-heat-transport:4}
(a)~$2.82 \times 365.25 = \qty{1030}{mm}$, \qty{1.03}{m} al año. (b)~$1.030 \times
78.1 = \qty{80.4}{W.m^{-2}}$.
(c)~El $LE$ del reanálisis es \qty{2}{W.m^{-2}} mayor. A lo largo de muchos años la atmósfera no puede seguir acumulando vapor de agua: toda el agua evaporada vuelve a caer, así que $LE$ y $LP$ deben ser iguales en la media global. La diferencia es un desequilibrio del modelo, como su desequilibrio radiativo.
\end{solution}

\begin{solution}{exo:b1:radiation-heat-transport:5}
(a)~$\bar Q = S_0(a/r)^2\sin\delta = 1361 \times 0.9683 \times 0.3978 =
\qty{524.2}{W.m^{-2}}$. (b)~$1361 \times 1.0334 \times 0.3978 =
\qty{559.5}{W.m^{-2}}$; el cociente es $1.0334/0.9683 = 1.067$. (c)~El polo tiene
luz solar 24 horas de 24 a una altura constante, así que su $\sin h_\odot$ medio a lo largo del día es $\sin\delta = 0.398$. En el ecuador el Sol está alto al mediodía pero bajo la mitad del tiempo; la media de $\sin h_\odot$ a lo largo de las 24 horas es $\cos\delta/\pi =
0.292$. El cociente es $\pi\tan\delta = 1.362$
(\cref{prop:b1:radiation-heat-transport:polar-summer}).
\end{solution}

\begin{solution}{exo:b1:radiation-heat-transport:6}
(a)~$\cos h_0 = -\tan\ang{40}\tan\ang{23.44} = -0.3638$, $h_0 = \ang{111.33} =
\qty{1.943}{rad}$; el día dura $24 \times 111.33/180 = \qty{14.84}{h}$; $\bar Q =
(1361/\pi) \times 0.9683 \times (1.943 \times 0.6428 \times 0.3978 + 0.7660 \times
0.9175 \times 0.9315) = \qty{483.0}{W.m^{-2}}$. (b)~$\cos h_0 = +0.3638$, $h_0 =
\ang{68.67}$, \qty{9.16}{h}; $\bar Q = \qty{155.9}{W.m^{-2}}$. (c)~El cociente es
$483.0/155.9 = 3.10$. La distancia contribuye $0.9683/1.0334 = 0.937$: trabaja en contra del verano boreal. La geometría, la altura del Sol y la duración del día contribuyen $3.10/0.937 = 3.31$, todo el contraste estacional y más.
\end{solution}

\begin{solution}{exo:b1:radiation-heat-transport:7}
(a)~$B = H/LE$: \num{0.093} en el ecuador, \num{0.374} a \ang{25} N,
\num{0.338} a \ang{45} N, \num{-0.152} a \ang{70} N. (b)~\ang{25} N: la media zonal incluye los grandes desiertos cálidos, tierra seca con poca agua que evaporar, donde el suelo caliente calienta el aire en su lugar. (c)~$H < 0$ significa que el calor fluye del aire a la superficie: en la media anual la nieve y el hielo marino están más fríos que el aire de encima, que se calienta con aire traído de \omterm{def:b1:spherical-astronomy:geographic}{latitudes} más bajas.
\end{solution}

\begin{solution}{exo:b1:radiation-heat-transport:8}
(a)~Evaluado: $\alpha = 100/340 = 0.294$, $R = 340 - 100 - 239 =
\qty{+1}{W.m^{-2}}$ (el \num{0.7} redondeado). Reanálisis: $\alpha = 117.0/341.9 =
0.342$, $R = 341.9 - 117.0 - 237.7 = \qty{-12.8}{W.m^{-2}}$ con estos términos redondeados,
\qty{-12.9}{W.m^{-2}} con los no redondeados de
\cref{ex:b1:radiation-heat-transport:global}. (b)~El reanálisis
absorbe $224.8$ frente a $240$, \qty{15.2}{W.m^{-2}} menos, porque refleja
\qty{17}{W.m^{-2}} más; su OLR es solo \qty{1.3}{W.m^{-2}} menor. El exceso de
reflexión explica el desequilibrio. (c)~Integrando $R$ desde el polo sur hasta el
polo norte se obtiene la integral global, $-12.9 \times 4\pi a_\oplus^2 = -12.9 \times
5.10\times10^{14} = \qty{-6.6}{PW}$. Un transporte a través del polo norte es
imposible: no hay nada más allá, el paralelo es un punto. La proposición
solo vale si $\int R\,\mathrm dA = 0$, lo que exige el estado estacionario.
\end{solution}

\begin{solution}{exo:b1:radiation-heat-transport:9}
(a)~El área es $2\pi a_\oplus^2 \times 2\sin\ang{30} = 2\pi a_\oplus^2 =
\qty{2.55e14}{m^2}$, la mitad de la Tierra; el excedente es $33 \times 2.55\times10^{14} =
\qty{8.4}{PW}$. (b)~\qty{4.2}{PW} a través de cada paralelo. (c)~Los dos suman
\qty{8.5}{PW}, el excedente, pero son desiguales: el hemisferio sur exporta
\qty{0.5}{PW} más que el norte. El excedente no se reparte en el ecuador: en este reanálisis \qty{0.36}{PW} del excedente de los trópicos del norte cruza el ecuador hacia el sur.
\end{solution}

\begin{solution}{exo:b1:radiation-heat-transport:10}
(a)~Escríbase la media diaria como $\bar Q = (a/r)^2 F(\varphi, \delta)$, con $F$ el
corchete de \cref{thm:b1:radiation-heat-transport:daily} multiplicado por $S_0/\pi$. Durante
$\mathrm dt$ la energía recibida por metro cuadrado es $\bar Q\,\mathrm dt$, y
$\mathrm dt = r^2\,\mathrm d\lambda/h$, así que $\bar Q\,\mathrm dt = (a^2/h)
F(\varphi, \delta(\lambda))\,\mathrm d\lambda$: el factor $r^2$ se cancela.
(b)~A lo largo de un año la energía es $(a^2/h)\int_0^{2\pi} F(\varphi,
\delta(\lambda))\,\mathrm d\lambda$. La fórmula no cambia cuando $\varphi$ y
$\delta$ cambian de signo a la vez, $F(-\varphi, \delta) = F(\varphi, -\delta)$, y
$\delta(\lambda + \pi) = -\delta(\lambda)$, pues $\sin\delta = \sin\varepsilon
\sin\lambda$. Desplazar $\lambda$ en $\pi$ dentro de la integral muestra que $-\varphi$
recibe la misma energía que $\varphi$. (c)~La primavera y el verano boreales duran
\num{7.6} días más, porque la Tierra está cerca del afelio y se mueve despacio entonces; el verano austral es más brillante pero más corto. Por (a), los dos efectos se compensan exactamente.
\end{solution}

\begin{solution}{exo:b1:radiation-heat-transport:11}
(a)~$\langle (a/r)^2\rangle = \frac1P\int_0^P \frac{a^2}{r^2}\,\mathrm dt =
\frac1P\int_0^{2\pi}\frac{a^2}{r^2}\frac{r^2}{h}\,\mathrm d\lambda = \frac{2\pi
a^2}{hP} = \frac{2\pi a^2}{2\pi a b} = \frac ab$, y $b = a\sqrt{1-e^2}$.
(b)~La media global instantánea $\frac14 S_0(a/r)^2$ promediada da $\frac14 S_0
\langle(a/r)^2\rangle = S_0/(4\sqrt{1-e^2})$. (c)~$340.25/\sqrt{1 - 0.0167^2} =
\qty{340.30}{W.m^{-2}}$, \qty{0.05}{W.m^{-2}} o \qty{0.014}{\percent} por encima de
$S_0/4$.
\end{solution}

\begin{solution}{exo:b1:radiation-heat-transport:12}
(a)~Atmósfera $0.78 \times 4.10 = \qty{3.20}{PW}$, océano $0.22 \times 4.10 =
\qty{0.90}{PW}$. (b)~Se encuentra que la atmósfera sola transporta más, \qty{5.0}{PW}, que la atmósfera y el océano juntos del reanálisis, \qty{4.02}{PW}: el transporte total del reanálisis se queda corto al menos en una quinta parte ($1 - 4.02/5.0 = 0.20$).
Su \omterm{def:b1:radiation-heat-transport:net-radiation}{radiación neta} varía demasiado poco entre los trópicos y los polos. (c)~Quitar una constante de $R$ hace nula su integral global, pero deja sin cambio su variación con la \omterm{def:b1:spherical-astronomy:geographic}{latitud}; el error está en esa variación, en el albedo de las nubes y de las superficies del modelo en cada \omterm{def:b1:spherical-astronomy:geographic}{latitud}, no solo en la media global.
\end{solution}

\begin{solution}{pb:b1:radiation-heat-transport:1}
\begin{enumerate}
  \item Junio: $\cos h_0 = -0.3638$, $h_0 = \ang{111.33}$, $\bar Q =
        \qty{483.0}{W.m^{-2}}$. Diciembre: $h_0 = \ang{68.67}$, $\bar Q =
        \qty{155.9}{W.m^{-2}}$. Marzo: $h_0 = \ang{90}$, $\bar Q = (1361/\pi)
        \times 1.0081 \times \cos\ang{40} = \qty{334.6}{W.m^{-2}}$.
  \item Junio: $\tan\ang{80}\tan\ang{23.44} = 2.46 > 1$, día polar, $h_0 = \pi$:
        $\bar Q = 1361 \times 0.9683 \times \sin\ang{80}\sin\ang{23.44} =
        \qty{516.3}{W.m^{-2}}$. Diciembre: noche polar, $\bar Q = 0$. Marzo:
        $(1361/\pi) \times 1.0081 \times \cos\ang{80} = \qty{75.8}{W.m^{-2}}$.
  \item A \ang{80} N el Sol no se pone ($h_0 = \ang{180}$, 24 horas), mientras que a
        \ang{40} N brilla \qty{14.8}{h}. Su altura es menor, entre \ang{13.4} y \ang{33.4}, pero el término $h_0\sin\varphi\sin\delta$, que crece con la duración del día y la \omterm{def:b1:spherical-astronomy:geographic}{latitud}, gana.
  \item Los valores de las bandas superan a los valores centrales en \qty{0.5}{\percent},
        \qty{0.8}{\percent}, \qty{1.7}{\percent}, \qty{1.7}{\percent} y
        \qty{1.3}{\percent}. Primero, el reanálisis usa una irradiancia solar algo mayor (pregunta~5), un \qty{0.5}{\percent} más. Segundo, la media de una banda no es el valor en el centro: la mitad de cada banda más cercana al ecuador tiene más área y más luz solar que la mitad más cercana al polo.
  \item Por \cref{prop:b1:radiation-heat-transport:global}, $S_0 = 4\sqrt{1-e^2}
        \times 341.9 = \qty{1367}{W.m^{-2}}$, \qty{0.47}{\percent} más que \qty{1361}{W.m^{-2}}.
  \item $A = 2\pi a_\oplus^2(1 - \sin\ang{40}) = 2\pi \times (6.371\times10^6)^2
        \times 0.3572 = \qty{9.110e13}{m^2}$, una fracción $(1 - \sin\ang{40})/2 =
        0.179$ de la Tierra.
  \item $(1 - \alpha)\bar Q$: \num{198.0}, \num{159.3}, \num{124.7}, \num{92.5},
        \num{74.7} \unit{W.m^{-2}}.
  \item $R$: \num{-28.1}, \num{-51.0}, \num{-76.6}, \num{-104.9}, \num{-116.9}
        \unit{W.m^{-2}}. Las cinco bandas están en déficit; la banda polar
        \ang{80}--\ang{90} es la mayor.
  \item $T_e = (\text{OLR}/\sigma)^{1/4}$: \num{251.3}, \num{246.8}, \num{244.1},
        \num{242.9}, \num{241.1} \unit{K}.
  \item Absorbida: $198.0/74.7 = 2.65$; OLR: $226.1/191.6 = 1.18$.
  \item El Sol está más bajo, y la luz que llega con un ángulo rasante es reflejada más por las nubes y el agua; la nieve y el hielo marino, superficies brillantes, cubren más suelo y mar hacia el polo.
  \item La OLR va como $T_e^4$, y $T_e$ varía solo \qty{10}{K}, \qty{4}{\percent},
        a lo largo del casquete: las altas \omterm{def:b1:spherical-astronomy:geographic}{latitudes} se mantienen mucho más cálidas de lo que permitiría su luz solar sola, por el calor traído desde el sur.
  \item En \unit{m^2}: \num{3.143e13}, \num{2.550e13}, \num{1.879e13},
        \num{1.151e13}, \num{0.387e13}; la suma es \qty{9.110e13}{m^2}.
  \item $R_iA_i$: \num{-0.883}, \num{-1.300}, \num{-1.439}, \num{-1.207},
        \num{-0.453} \unit{PW}; el casquete: \qty{-5.28}{PW}.
  \item $\alpha_\text{casquete} = \sum\alpha_i\bar Q_iA_i/\sum\bar Q_iA_i = 0.400$,
        frente a \num{0.342} del globo.
  \item $-12.90 \times 5.10\times10^{14} = \qty{-6.58}{PW}$: la Tierra perdería calor a ese ritmo y se enfriaría. El desequilibrio evaluado es
        $+0.7 \times 5.10\times10^{14} = \qty{+0.36}{PW}$, dieciocho veces menor y
        de signo opuesto: el desequilibrio del reanálisis es un error del modelo.
  \item Corregido $R_i + 12.90$: \num{-15.2}, \num{-38.1}, \num{-63.7},
        \num{-92.0}, \num{-104.0} \unit{W.m^{-2}}; potencias \num{-0.478},
        \num{-0.972}, \num{-1.197}, \num{-1.059}, \num{-0.403} \unit{PW}; el casquete:
        \qty{-4.108}{PW}, un déficit medio de \qty{45.1}{W.m^{-2}}.
  \item La energía del casquete no cambia de un año a otro, así que lo que pierde en el tope de la atmósfera debe entrar en él, y la única vía de entrada es a través del paralelo 40: $T(\ang{40}) = -2\pi a_\oplus^2\int_{\ang{40}}^{\ang{90}}
        R\cos\varphi\,\mathrm d\varphi$.
  \item $T(\ang{40}\,\text{N}) = +\qty{4.11}{PW}$, hacia el norte.
  \item $LP = 2.4655\times10^6 \times 1.95/86400 = \qty{55.6}{W.m^{-2}}$, más que
        $LE = \qty{42.4}{W.m^{-2}}$: cae más agua sobre el casquete de la que se evapora de él; la diferencia llega como vapor.
  \item $(LE - LP) - 2.07 = 42.4 - 55.6 - 2.07 = \qty{-15.3}{W.m^{-2}}$; multiplicado por el
        área del casquete: \qty{-1.40}{PW}. La atmósfera lleva \qty{1.40}{PW}
        de calor latente a través de \ang{40} N.
  \item $1.40/4.11 = 0.34$: un tercio.
  \item Océano $0.22 \times 4.11 = \qty{0.90}{PW}$; atmósfera $0.78 \times 4.11 =
        \qty{3.20}{PW}$, de los cuales \qty{1.40}{PW} como calor latente y \qty{1.80}{PW}
        transportados como calor sensible y energía potencial del aire.
  \item Una cota inferior. El total del reanálisis a \ang{43} N, \qty{4.02}{PW}, es menor que la atmósfera sola en el estudio, \qty{5.0}{PW}; el total verdadero es mayor.
  \item Alrededor de \qty{4.11}{PW}, \qty{4.11e15}{W}, según la radiación del reanálisis, un tercio de ello como calor latente; el valor verdadero es aún mayor.
\end{enumerate}
\end{solution}
